% 本文件由 LaTeX 排版 Agent 根据有效 translation.json 自动生成。
% author=Gregory Thaumaturgus; work_id=0606; title=论天主圣三残篇

\chapter{论天主圣三残篇}

% structure: p0001 source=h1 target=omitted reason=page-title-already-rendered-by-parent

显灵迹者格里高利，庞图斯新凯撒利亚的主教，接近宗徒的继承者，在其论天主圣三的讲辞中如此说：—\par

我在三者中都看到要素：实体、类别、名称。我们说到人、仆人、\OriginalTerm{管理人}{curatorem}——人就其实体而言；仆人就其类别或状况而言；管理人则就其称谓而言。我们也说到圣父、圣子、圣神；然而，这些并不是后来才附加的称谓，而是\OriginalTerm{自立体}{subsistences}。再者，\InnerQuote{人}这一称谓并非实际上只是一个称谓，而是众人共有的实体，是适合所有人的称谓。此外，名字是这样的：\Person{亚当}、\Person{亚伯拉罕}、\Person{依撒格}、\Person{雅各伯}：我说，这些是名字。但神圣的位格确实是名字：而这些名字就是位格；位格意指那存在且自立者——即天主的本质。本性的名称也意指自立体；就如我们说到\InnerQuote{人}。所有的（位格）是同一性体、同一本质、同一意志，被称为圣三；这些也都是自立的名称，三位同一性体，同一类别。但圣子的位格在其一体性中（\OriginalTerm{一体性}{unita est}）是复合的，是由二者——即神性与人性结合——而成为一的，二者构成一体。然而神性并不因而增加任何成分，圣三依然如故。亦无任何新事临于位格或名称，而是永恒且超越时间。然而，无人足以知道这些，直到圣子降生成人而启示之，说：\BibleQuote{父啊！我已将祢的名彰显给世人；求祢也光荣我，使他们知道我是祢的子。} \BibleRef{若 17:6} 在山上圣父发言说：\BibleQuote{这是我的爱子。} \BibleRef{玛 3:17} 同一位（圣父）在约旦河派遣了祂的圣神。如此便向我们宣告了存在一个同等尊荣的永恒圣三。此外，圣子由圣父所生是不可思议、不可言喻的；因为它是属神的，追究它便成为不可能：因为属神的事物不能被属肉体的事物理解或追寻，因它远超人性。我们人类的确知道我们自身的生育，以及其它物体的生育；但属神之事超越人类处境，丝毫不能被人的心灵理解。属神的实体不能朽坏也不能分解；然而我们的，如容易了解的，会朽坏、会分解。诚然，人如何可能——他受限于六面：东、西、南、北、深渊、天空——理解那超越诸天、在深渊之下、延伸至南北之外、临在每一处、充满一切虚空的事物呢？但如果我们真能审察属神的实体，它的卓越性倒真会被破坏。让我们思考在我们身体内所进行的事情；并且进一步，让我们看看我们是否有能力确知意念如何从心生出，言语如何从舌生出，以及诸如此类。如今，如果我们丝毫不能领会自身内所行的事，我们怎能理解那超越一切心智的未受造造物主的奥秘呢？确实，如果这奥秘是人能穿透的，受默感的\Person{若望}就绝不会如此断言：\BibleQuote{从来没有人见过天主。} \BibleRef{若 1:18} 那么，那从来没有人见过的一位——我们能拿什么来类比祂，好借此理解祂的受生呢？我们确实毫不含糊地知道我们的灵魂内在于我们，与身体结合；但谁曾见过自己的灵魂？谁曾辨别出它与身体的结合？我们确切知道的仅此一事：我们体内有一个灵魂与身体结合。如此，我们推论并相信圣言是由圣父所生，尽管我们不拥有也不知晓这事实清晰的\OriginalTerm{理由}{rationale}。圣言本身在一切受造物之前——从永恒者而来的永恒者，如泉水从泉水，光明从光明。\InnerQuote{圣言}这一语词，确实属于圣经中提及的那三类非实质性的言语——即\Emph{默想之言语}（\OriginalTerm{默想之言语}{conceived}）、\Emph{说出之言语}（\OriginalTerm{说出之言语}{uttered}）和\Emph{发音之言语}（\OriginalTerm{发音之言语}{articulated}）。默想之言语当然不是实质性的。说出之言语则是先知从天主听到的声音，或先知的话语本身；这也不是实质性的。最后，发音之言语是由词句组成的、在空气中形成的人的话语（\OriginalTerm{在空气中形成}{aere efformatus}），这也不是实质性的。但天主的圣言是实质性的，具有崇高而常存的性体，与祂同永恒，与祂不可分离，永不失落，而永久处于永恒结合中。这圣言创造了天地，万有都是在祂内受造的。祂是天主的手臂和能力，因不可分的性体，从未与圣父分离，与圣父一起无始。这圣言从童贞玛利亚取了我们的实体；就祂是属神而言，祂与圣父不可分地相等；就祂是属肉体而言，祂同样与我们不可分离地相等。并且，就祂是属神而言，祂以\OriginalTerm{同等}{aequiparat}不失、无限地赐予圣神。在圣言——圣子——降生成人之前，并非有两个性体，而是只有圣三的一个性体；在圣子降生成人之后，圣三的性体仍然是一个。但若有人因圣言摄取人性而相信圣三有所增加，他便与我们以及天主教会与宗徒教会的职务不相干。这是神圣天主圆满、神圣、宗徒的信德。赞美圣三，永无穷尽，世世无尽。阿们。\par

\SourceRecord{Translated by S.D.F. Salmond.；Ante-Nicene Fathers；Vol. 6.；Edited by Alexander Roberts, James Donaldson, and A. Cleveland Coxe.；Buffalo, NY: Christian Literature Publishing Co.,；1886.；<http://www.newadvent.org/fathers/0606.htm>.}
